\subsection{TENSOR PRODUCTS OF VECTOR SPACES}

The results of §§ 3, 4 and 5 relating to free or projective modules apply in particular to vector spaces and give the following properties:

PROPOSITION 14. Given an exact sequence

\[
0 \rightarrow \mathbf {E} ^ {\prime} \rightarrow \mathbf {E} \rightarrow \mathbf {E} ^ {\prime \prime} \rightarrow 0\tag{23}
\]

of right vector spaces over a field K and linear mappings and a left vector space F over K, the corresponding sequence of Z-linear mappings

\[
0 \rightarrow \mathrm{E} ^ {\prime} \otimes_ {\mathbb {K}} \mathrm{F} \rightarrow \mathrm{E} \otimes_ {\mathbb {K}} \mathrm{F} \rightarrow \mathrm{E} ^ {\prime \prime} \otimes_ {\mathbb {K}} \mathrm{F} \rightarrow 0
\]

is exact and splits.

As the sequence (23) splits, this is a particular case of § 3, no. 7, Corollary 5 to Proposition 7 and § 3, no. 6, Proposition 5.

Because of Proposition 14, when \(\mathbf{E}'\) is a vector subspace of \(\mathbf{E}\) , \(j: \mathbf{E}' \to \mathbf{E}\) the canonical injection, \(\mathbf{E}' \otimes_{\mathbb{K}} \mathbf{F}\) is usually identified with a sub- \(\mathbf{Z}\) -module of \(\mathbf{E} \otimes_{\mathbb{K}} \mathbf{F}\) by means of the injection \(j \otimes 1_{\mathbf{F}}\) . With this convention:

COROLLARY. Let K be a field, E a right vector space over K, F a left vector space over K, \((\mathbf{M}_{\alpha})_{\alpha \in \mathbf{A}}\) a family of vector subspaces of E and \((\mathbf{N}_{\beta})_{\beta \in \mathbf{B}}\) a family of vector subspaces of F. Then

\[
\left(\bigcap_ {\alpha \in \mathrm{A}} \mathrm{M} _ {\alpha}\right) \otimes_ {\mathrm{K}} \left(\bigcap_ {\beta \in \mathrm{B}} \mathrm{N} _ {\beta}\right) = \bigcap_ {(\alpha , \beta) \in \mathrm{A} \times \mathrm{B}} \left(\mathrm{M} _ {\alpha} \otimes_ {\mathrm{K}} \mathrm{N} _ {\beta}\right).\tag{24}
\]

It is obviously sufficient to prove the particular case

\[
\left(\bigcap_ {\alpha \in \mathrm{A}} \mathrm{M} _ {\alpha}\right) \otimes_ {\mathrm{K}} \mathrm{F} = \bigcap_ {\alpha \in \mathrm{A}} \left(\mathrm{M} _ {\alpha} \otimes_ {\mathrm{K}} \mathrm{F}\right).\tag{25}
\]

Clearly the left hand side of (25) is contained in the right hand side. To prove the converse, we consider a basis \((f_{\lambda})_{\lambda \in L}\) of F. Every element of \(\mathbf{E} \otimes_{\mathbb{K}} \mathbf{F}\) can then be expressed uniquely in the form \(\sum_{\lambda \in L} x_{\lambda} \otimes f_{\lambda}\) , where \(x_{\lambda} \in \mathbf{E}\) (§3, no. 7, Corollary 1 to Proposition 7); if \(\mathbf{E}'\) is a vector subspace of E, the relation \(\sum_{\lambda \in L} x_{\lambda} \otimes f_{\lambda} \in \mathbf{E}' \otimes_{\mathbb{K}} \mathbf{F}\) is equivalent, by Proposition 14, to \(x_{\lambda} \in \mathbf{E}'\) for all \(\lambda \in L\) . To say that \(\sum_{\lambda \in L} x_{\lambda} \otimes f_{\lambda}\) belongs to each of the \(\mathbf{M}_{\alpha} \otimes_{\mathbb{K}} \mathbf{F}\) thus means that for all \(\lambda \in L\) and all \(\alpha \in A\) , \(x_{\lambda} \in M_{\alpha}\) , that is \(x_{\lambda} \in \bigcap_{\alpha \in A} M_{\alpha}\) for all \(\lambda \in L\) , which proves that the right hand side of (25) is contained in the left hand side. PROPOSITION 15. If \((E_{\lambda})_{\lambda \in L}\) is a family of right vector spaces over a field K and \((F_{\mu})_{\mu \in M}\) a family of left vector spaces over K, the canonical mapping

\[
\left(\prod_ {\lambda \in L} E _ {\lambda}\right) \otimes_ {K} \left(\prod_ {\mu \in M} F _ {\mu}\right)\rightarrow \prod_ {(\lambda , \mu) \in L \times M} (E _ {\lambda} \otimes_ {K} F _ {\mu})\tag{26}
\]

(§ 3, no. 7, formula (22)) is injective.

We write \(\mathbf{F} = \prod_{\mu \in \mathbf{M}} \mathbf{F}_{\mu}\) ; the mapping (26) is the composition of the canonical mappings

\[
\left(\prod_ {\lambda \in L} \mathbf {E} _ {\lambda}\right) \otimes_ {\mathbb {K}} F \rightarrow \prod_ {\lambda \in L} \left(\mathbf {E} _ {\lambda} \otimes_ {\mathbb {K}} F\right)
\]

and

\[
\prod_ {\lambda \in L} \left(\mathrm{E} _ {\lambda} \otimes_ {K} F\right)\rightarrow \prod_ {\lambda \in L} \left(\prod_ {\mu \in M} \left(\mathrm{E} _ {\lambda} \otimes_ {K} F _ {\mu}\right)\right);
\]

as F and the \(\mathbf{E}_{\lambda}\) are vector spaces over K, this reduces to §3, no. 7, Corollary 3 to Proposition 7.

When the conditions of Proposition 15 are fulfilled, the tensor product \(\left(\prod_{\lambda \in L} E_{\lambda}\right) \otimes_{K} \left(\prod_{\mu \in M} F_{\mu}\right)\) is often identified with its canonical image in

\[
\prod_ {\lambda , \mu} (\mathbf {E} _ {\lambda} \otimes_ {K} \mathbf {F} _ {\mu}).
\]

With this convention:

COROLLARY. Let F be a left vector space over K; for every set X, the left vector space \(K_{d}^{X} \otimes_{K} F\) is identified with the subspace of the space \(F^{X}\) of all mappings of X into F, consisting of the mappings u such that \(u(X)\) is of finite rank in F.

If \((f_{\lambda})\) is a basis of \(F\) , the element \(\sum_{\lambda \in L} v_{\lambda} \otimes f_{\lambda}\) of \(K_d^X \otimes_K F\) is identified under (26) with the mapping \(x \mapsto \sum_{\lambda} v_{\lambda}(x)f_{\lambda}\) . As \(v_{\lambda} = 0\) except for the indices \(\lambda\) belonging to a finite subset \(H\) of \(L\) , the image of \(X\) under the above mapping is contained in the finite-dimensional subspace of \(F\) generated by the \(f_{\lambda}\) of indices \(\lambda \in H\) . Conversely, let \(u: X \to F\) be a mapping such that \(u(X)\) is contained in a finite-dimensional subspace \(G\) of \(F\) and let \((b_i)_{1 \leqslant i \leqslant n}\) be a basis of \(G\) . For all \(x \in X\) , we may write \(u(x) = \sum_{i=1}^{n} v_i(x)b_i\) , where the \(v_i(x)\) are well determined elements of \(K\) ; thus \(n\) mappings \(v_i: X \to K\) are defined and clearly \(u\) is then identified with the element \(\sum_{i=1}^{n} v_i \otimes b_i\) .

Similarly, for a right vector space E over K and a set Y, \(E \otimes_{K} K_{s}^{Y}\) is identified with a subspace of the space \(E^{Y}\) , consisting of the mappings \(v: Y \to E\) such that \(v(Y)\) is of finite rank. More particularly, for every field K, \(K_{d}^{X} \otimes_{K} K_{s}^{Y}\) is identified with a subspace of the space \(K^{X \times Y}\) of mappings of \(X \times Y\) into K (K being considered as a (K, K)-bimodule); an element \(\sum_{i} u_{i} \otimes v_{i}\) , where \(u_{i}\) is a mapping of X into K and \(v_{i}\) a mapping of Y into K, is identified with the mapping \((x, y) \mapsto \sum_{i} u_{i}(x) v_{i}(y)\) of \(X \times Y\) into K.

PROPOSITION 16. (i) Let K, L be two fields, E a left vector space over K, F a left vector space over L and G a (K, L)-bimodule. Then the canonical Z-homomorphism

\[
\nu : \operatorname{Hom} _ {\mathbf {K}} (\mathrm{E}, \mathrm{G}) \otimes_ {\mathrm{L}} \mathrm{F} \rightarrow \operatorname{Hom} _ {\mathbf {K}} (\mathrm{E}, \mathrm{G} \otimes_ {\mathrm{L}} \mathrm{F})\tag{27}
\]

(§ 4, no. 2, formula (7)) is injective; it is bijective when one of the vector spaces E, F is finite-dimensional.

\begin{enumerate}
\def\labelenumi{(\roman{enumi})}
\setcounter{enumi}{1}
\tightlist
\item
  Let \(E_{1}\) , \(E_{2}\) , \(F_{1}\) , \(F_{2}\) be four vector spaces over a commutative field K; then the canonical K-homomorphism
\end{enumerate}

\[
\lambda : \operatorname{Hom} \left(\mathrm{E} _ {1}, \mathrm{F} _ {1}\right) \otimes \operatorname{Hom} \left(\mathrm{E} _ {2}, \mathrm{F} _ {2}\right)\rightarrow \operatorname{Hom} \left(\mathrm{E} _ {1} \otimes \mathrm{E} _ {2}, \mathrm{F} _ {1} \otimes \mathrm{F} _ {2}\right)\tag{28}
\]

(§ 4, no. 4, formula (21)) is injective; it is bijective if one of the ordered pairs \((\mathbf{E}_1, \mathbf{E}_2), (\mathbf{E}_1, \mathbf{F}_1), (\mathbf{E}_2, \mathbf{F}_2)\) consists of finite-dimensional spaces.

Assertion (i) is a particular case of § 4, no. 2, Proposition 2. Similarly the second assertion of (ii) is a particular case of § 4, no. 4, Proposition 4. Finally, to see that the homomorphism (28) is always injective, observe that \(\operatorname{Hom}(\mathbf{E}_i, \mathbf{F}_i)\) is a vector subspace of \(\mathbf{F}_i^{\mathbf{E}_i}\) ( \(i = 1, 2\) ) and that

\[
\operatorname{Hom} \left(\mathrm{E} _ {1} \otimes \mathrm{E} _ {2}, \mathrm{F} _ {1} \otimes \mathrm{F} _ {2}\right)
\]

is canonically identified with a vector subspace of the space \((\mathbf{F}_{1} \otimes \mathbf{F}_{2})^{\mathbf{E}_{1} \times \mathbf{E}_{2}}\) (II, §3, no. 1, Proposition 1); when these identifications are made and also the left hand side of (28) is identified with a subspace of \(F_{1}^{E_{1}} \otimes F_{2}^{E_{2}}\) (Proposition 14), the canonical mapping (28) becomes the restriction to this subspace of the canonical mapping (26) and it has been seen (Proposition 15) that the latter is injective.

COROLLARY 1. Let E and F be two vector spaces over a field K; the canonical mapping

\[
\mathbf {E} ^ {*} \otimes_ {\mathbb {K}} \mathbf {F} \rightarrow \operatorname{Hom} _ {\mathbb {K}} (\mathbf {E}, \mathbf {F})
\]

(§ 4, no. 2, formula (11)) is injective; it is bijective when E or F is finite-dimensional.

This is a special case of Proposition 16 (i).

COROLLARY 2. Let E be a right vector space and F a left vector space over the same field K; the canonical mapping

\[
\mathbf {E} \otimes_ {\mathbb {K}} \mathbf {F} \rightarrow \operatorname{Hom} _ {\mathbb {K}} (\mathbf {E} ^ {*}, \mathbf {F})\tag{29}
\]

(§ 4, no. 2, formula (15)) is injective; it is bijective when E is finite-dimensional.

This is a special case of § 4, no. 2, Remark 2.

Remarks. (1) Let K be a commutative field, E, F two vector spaces over K,

\((a_{\lambda})\) a basis of \(\mathbf{E}\) and \((b_{\mu})\) a basis of \(\mathrm{F}\) ; then \((a_{\lambda} \otimes b_{\mu})\) is a basis of the vector \(\mathbf{K}\) -space \(\mathbf{E} \otimes_{\mathbb{K}} \mathbf{F}\) (II, § 3, no. 7, Corollary 2 to Proposition 2) and therefore

\[
\dim_ {\mathbf {K}} (\mathbf {E} \otimes_ {\mathbf {K}} \mathbf {F}) = \dim_ {\mathbf {K}} \mathbf {E}. \dim_ {\mathbf {K}} \mathbf {F}.\tag{30}
\]

\begin{enumerate}
\def\labelenumi{(\arabic{enumi})}
\setcounter{enumi}{1}
\tightlist
\item
  Let K be a commutative field, \(E_{1}\) , \(E_{2}\) , \(F_{1}\) , \(F_{2}\) four vector spaces over K and \(u: E_{1} \rightarrow F_{1}\) , \(v: E_{2} \rightarrow F_{2}\) two linear mappings; then
\end{enumerate}

\[
\operatorname{rg} (u \otimes v) = \operatorname{rg} (u). \operatorname{rg} (v).\tag{31}
\]

It is immediate that \((u\otimes v)(\mathbf{E}_1\otimes \mathbf{E}_2)\) is the canonical image of \(u(\mathbf{E}_1)\otimes v(\mathbf{E}_2)\) in \(\mathbf{F}_1\otimes \mathbf{F}_2\) and hence (Proposition 14) is isomorphic to \(u(\mathbf{E}_1)\otimes v(\mathbf{E}_2)\) ; the conclusion then follows from (30).

\begin{enumerate}
\def\labelenumi{(\arabic{enumi})}
\setcounter{enumi}{2}
\tightlist
\item
  Under the same hypotheses as in Remark 1,
\end{enumerate}

\[
\dim_ {\mathbb {K}} (\operatorname{Hom} _ {\mathbb {K}} (\mathrm{E}, \mathrm{F})) \geqslant \dim_ {\mathbb {K}} \mathrm{E}. \dim_ {\mathbb {K}} \mathrm{F}.\tag{32}
\]

If E is isomorphic to \(K^{(I)}\) , \(\operatorname{Hom}(E, F)\) is isomorphic to \((\operatorname{Hom}(K, F))^{I}\) (§ 1, no. 6, Corollary 1 to Proposition 6) and hence to \(F^{I}\) (§ 1, no. 14); as \(F^{(I)}\) is a subspace of \(F^{I}\) and \(\dim(F^{(I)}) = \operatorname{Card}(I).dim F = \dim E. dim F\) (no. 2, Proposition 2), the inequality (32) follows from no. 3, Corollary 4 to Proposition 4. The same argument shows that the two sides of (32) are equal when dim E is finite (cf.~§ 10, nos. 3 and 4).

